\section{Rational inverse chains in a height-two formal group}\label{sec:cyclic}
This section proves the cyclic-point restriction by following the Newton slopes at every inverse step.

\begin{lemma}\label{lem:inverse-polygons}
Let $R$ and $F$ satisfy the hypotheses of Theorem~\ref{thm:cyclic-bound}. There is an integral coordinate and a deformation parameter $u\in\mathfrak m_R$ such that, writing $a=v_R(p)$, $b=v_R(u)$, and $q=p^2$, the nonzero $p$-torsion valuations are as follows. If $b\ge ap/(p+1)$, all have valuation $a/(q-1)$. Otherwise they have valuations
\begin{equation}\label{eq:alpha-beta}
 \alpha=\frac{a-b}{p-1}\quad(p-1\text{ points}),\qquad
 \beta=\frac{b}{p(p-1)}\quad(q-p\text{ points}).
\end{equation}
Here $b=+\infty$ is permitted when $u=0$.

If $z$ is a nonzero torsion point of valuation $0<Z<a$, its $p$-preimages have valuations
\begin{equation}\label{eq:inverse-polygons}
 \begin{cases}
 Z/q & b\ge Z(p-1)/p,\\[2pt]
 (Z-b)/p\text{ with multiplicity }p,
 \quad\beta\text{ with multiplicity }q-p
 & b<Z(p-1)/p.
 \end{cases}
\end{equation}
The equality case belongs to the first line. Once an inverse chain takes a step of the first kind, or a step to $\beta$, all its later inverse steps divide the valuation by $q$.
\end{lemma}

\begin{proof}
The integral universal height-two coordinate has
\[
 [p](T)=pT+(1-p^{p-1})uT^p+\sum_{j>p}c_jT^j,
 \qquad c_j\in(p,u)R\ (p<j<q),\quad c_q\in R^\times.
\]
All coefficients are integral and degrees two through $p-1$ vanish. This is the specialization of \citet[Proposition~5.7, equations~(12.3)--(12.5), and Proposition~12.10]{HopkinsGross1994}; the complete local base may have any residue-field extension. The formula for the coefficient of $T^p$ also follows directly by comparing $f([p]T)=pf(T)$ in the normalized logarithm $f(T)=T+(u/p)T^p+O(T^{p^2})$.

The only possible lower Newton vertices of $[p](T)$ at positive roots are $(1,a),(p,b),(q,0)$. Indeed intermediate coefficients have valuation at least $\min(a,b)$, later coefficients are integral, and all the omitted points lie on or above the indicated negative-slope segments. The middle point is below the segment joining the other two exactly when $b<ap/(p+1)$. This proves~\eqref{eq:alpha-beta}.

For $[p](T)-z$, the point $(1,a)$ lies strictly above the lower hull because $Z<a$. The possible vertices are $(0,Z),(p,b),(q,0)$. The middle point is below their single chord exactly when $b<Z(p-1)/p$, giving~\eqref{eq:inverse-polygons}. Every inverse valuation is smaller than its input. If the first inequality holds at $Z$, it continues to hold at $Z/q$. If the next valuation is $\beta$, then
$b=p(p-1)\beta>\beta(p-1)/p$, so the same conclusion applies. Finally every nonzero first-level valuation is at most $a/(p-1)<a$, so the inverse formulas apply at every subsequent step.
\end{proof}

\begin{proof}[Proof of Theorem~\ref{thm:cyclic-bound}]
Let $P$ be rational of order $p^n$. All points in its successive multiplication chain are rational, and hence their positive coordinate valuations are integers. Put $N=p^2-1$ and $q=p^2$.

In the one-slope case the terminal valuation is $a/(Nq^{n-1})$. Integrality implies $a\ge Nq^{n-1}$, which is stronger than the desired inequality. In the two-slope case, if the first-level point $p^{n-1}P$ has valuation $\beta$, then all $n-1$ inverse steps are uniform. Thus $b\ge p(p-1)q^{n-1}$, and the strict two-slope inequality gives
\[
 a>\frac{p+1}{p}b\ge Nq^{n-1}.
\]
We may therefore assume that the first-level point has valuation $Z_0=\alpha$.

Call a step \emph{high} when it uses $(Z-b)/p$ in the second line of~\eqref{eq:inverse-polygons}. After $j$ consecutive high steps, the recurrence gives
\begin{equation}\label{eq:high-recurrence}
 Z_j=\frac{a/p^j-b}{p-1},\qquad
 a=p^j\bigl(b+(p-1)Z_j\bigr).
\end{equation}
Since $Z_j,b$ are integers, $j\le v_p(a)=v$.

Suppose the first non-high step is uniform, following $j\le\min(v,n-2)$ high steps. There remain $r=n-1-j\ge1$ uniform steps. Their terminal valuation is integral, so $Z_j\ge q^r$. The uniform-step inequality says $b\ge Z_j(p-1)/p$. Hence
\[
 a=p^j\bigl(b+(p-1)Z_j\bigr)
   \ge Np^{j-1}Z_j
   \ge Np^{2n-3-j}
   \ge Np^{2n-3-v}.
\]
If $j<v$, the last inequality is strict. If $j=v$ and equality held throughout, then $a=Np^{2n-3-v}$, and therefore $v_p(a)=2n-3-v>v$, since $v=j\le n-2$. This is a contradiction. Thus the final inequality is strict also at the Newton boundary.

Suppose instead that the first non-high step is low, following $j$ high steps. After it, $n-2-j$ uniform steps remain. Consequently
\[
 b\ge p(p-1)q^{n-2-j}.
\]
The strict high-versus-low condition at $Z_j$ is equivalent, using~\eqref{eq:high-recurrence}, to $a/p^j>(p+1)b$. Therefore
\[
 a>Np^{2n-3-j}\ge Np^{2n-3-v}.
\]
Finally, if all $n-1$ inverse steps are high, then $v\ge n-1$, $b\ge1$ and $Z_{n-1}\ge1$. Equation~\eqref{eq:high-recurrence} gives
\[
 a\ge p^n>Np^{n-2}\ge Np^{2n-3-v}.
\]
These cases exhaust every inverse chain and prove the theorem.
\end{proof}

\begin{remark}
For $p\nmid a$ this says $a>(p^2-1)p^{2n-3}$. The dependence on $v_p(a)$ comes entirely from the integrality assertion in~\eqref{eq:high-recurrence}. The Newton slopes are the familiar canonical-subgroup slopes; the inequality is a convenient consequence of their full inverse-chain classification.
\end{remark}

\begin{proposition}\label{prop:mixed-obstruction}
Let $p\ge5$, let $e=e_K$ and $v=v_p(e)$. Let $M/K$ be a height-two Barsotti--Tate group with potentially supersingular good reduction. Let $H/K$ have a two-step filtration by potentially good height-one groups whose dimensions sum to one. If
$M[p^n]\simeq H[p^n]$ for $n\ge2$, then
\begin{equation}\label{eq:mixed-obstruction}
 e>p^{2n-3-v}.
\end{equation}
\end{proposition}

\begin{proof}
Pass to the completed maximal unramified extension, which preserves $e$ and the proposed torsion comparison. Over this base a potentially good height-one character is either $\epsilon$ or $\chi_p\epsilon$, with $\epsilon$ finite of order dividing $p-1$. The determinant of a potentially supersingular height-two group is $\chi_p\epsilon_0$ with the same kind of finite character.

Transport the first height-one piece of $H[p^n]$ to a primitive cyclic subgroup $C$ of $M[p^n]$. If its character is $\epsilon$, twisting by $\epsilon^{-1}$ gives a rational point of exact order $p^n$. If its character is $\chi_p\epsilon$, first quotient $M$ by $C$. The image of $M[p^n]$ in the quotient is cyclic of order $p^n$, with character $\epsilon_0\epsilon^{-1}$; a finite twist again makes it rational. Schematic closure of $C$ in a good model is finite flat, so the quotient still acquires a good formal model of height two. Finite twisting preserves potential good reduction.

The faithful finite special inertia of a height-two formal group embeds in the units of the quaternion division algebra of invariant $1/2$ over $\Q_p$. Every element of this algebra satisfies a quadratic reduced polynomial. An element of order $p$ would generate $\Q_p(\zeta_p)$ of degree $p-1>2$, which is impossible. Its finite subgroups therefore have order prime to $p$ and inject, by reduction of the maximal order, into $\F_{p^2}^\times$. Thus a common faithful good-reduction field has tame degree $m\mid p^2-1$. The passage from faithful special inertia to a good model for the original lattice follows from crystalline descent and the weight-$\{0,1\}$ realization theorem of \citet[Corollary~2.2.6]{Kisin2006}; a finite-flat quotient adjusts the lattice if necessary.

Over that field the rational point is a formal point and the absolute index is $a=me$, with $v_p(a)=v$. Theorem~\ref{thm:cyclic-bound} gives $me>(p^2-1)p^{2n-3-v}$. Since $m\le p^2-1$, this implies~\eqref{eq:mixed-obstruction}.
\end{proof}
